\documentclass[10.5pt,a4paper]{article}
\usepackage[top=2.0cm,bottom=2.0cm,left=2.2cm,right=2.2cm]{geometry}
\usepackage{booktabs}
\usepackage{amsmath}
\usepackage{parskip}
\usepackage{enumitem}
\usepackage{titlesec}
\usepackage{microtype}
\usepackage{array}
\usepackage[hidelinks]{hyperref}

\setlength{\parskip}{4pt}
\titleformat{\section}{\normalsize\bfseries\uppercase}{}{0em}{}[\vspace{-2pt}\hrule\vspace{2pt}]
\titlespacing{\section}{0pt}{8pt}{4pt}
\titleformat{\subsection}{\normalsize\itshape}{}{0em}{}
\titlespacing{\subsection}{0pt}{5pt}{2pt}

\title{\textbf{Radar Feature Discriminability for Maritime Vessel Classification:\\
A Physics-Constrained Open Problem}\\[2pt]
\small\normalfont Prepared for expert consultation — strictly confidential,
not for circulation}
\author{\small Vijay Kumar V \quad|\quad Independent Researcher and Innovator,
Bangalore, India}
\date{\small August 2026}

\begin{document}
\maketitle\vspace{-8pt}
\hrule\vspace{6pt}

% ── Background ───────────────────────────────────────────────────────────────
\section{Background and Problem}

A land-based coastal surveillance radar collects amplitude returns from maritime
vessels at ranges of 5--50\,km. The goal is to classify vessel type
(Fishing / Tug / Cargo / Tanker) using only the radar return — with no reliance on
AIS transponder data, which can be spoofed or disabled. After exhaustive derivation
and statistical evaluation of every feature obtainable from the amplitude-only radar
output, it has been established that Cargo and Tanker vessels \textbf{cannot be
reliably discriminated} using the current feature set.

The conclusion is that this is a measurement physics constraint, not a processing
limitation. This document seeks expert guidance on whether the existing radar
receiver output contains additional physical information — accessible through
receiver tuning or processing parameter changes alone — that could provide the
needed discriminatory power.

\textbf{Hard constraint: no hardware modifications are possible.} The radar operates
within a secure facility. Transmitter, antenna, waveform, carrier frequency, and
polarisation are fixed. Only receiver settings, clutter filter parameters, integration
mode, and post-detection processing choices are adjustable.

% ── Dataset ──────────────────────────────────────────────────────────────────
\section{Dataset Summary}

\begin{center}\small
\begin{tabular}{lrrr}
\toprule
Vessel class & Detections & Tracks & Range (median) \\
\midrule
Fishing (small commercial fishing vessel) & 1,878 & 65 & 14\,km \\
Tug (harbour / ocean tug)                 & 1,823 & 37 & 20\,km \\
Cargo (bulk carrier / container ship)     & 4,500 & 264 & 30\,km \\
Tanker (oil / chemical / LNG tanker)      & 4,109 & 156 & 30\,km \\
\midrule
Total & 12,310 & 510 & \\
\bottomrule
\end{tabular}
\end{center}

Ground truth vessel type is from AIS records matched by position and time.
AIS-reported dimensions (length, beam) are excluded from all analysis: they are
transponder-provided labels, not radar observables.

% ── Features ─────────────────────────────────────────────────────────────────
\section{Features Derived from the Radar Return}

All features below are derived solely from the per-detection amplitude and spatial
extent outputs of the existing radar processing chain. The $R^{4}$ range correction
($\ln A + 4\ln R$) is applied to amplitude features to remove range-dependent
sensitivity variation.

\begin{center}\small
\begin{tabular}{ll}
\toprule
Feature & Physical basis \\
\midrule
Log peak RCS & Peak cell amplitude, $R^4$-corrected — vessel's dominant scatterer \\
Log total RCS & Sum of all cell amplitudes, $R^4$-corrected — total scattering power \\
RCS concentration & Peak / total ratio — point vs.\ distributed scatterer \\
Aspect ratio & Cross-range extent / down-range extent — vessel orientation proxy \\
Footprint area (m$^2$) & Product of range and azimuth extents — coarse vessel size \\
Down-range extent (m) & Detection span in range — vessel radial length proxy \\
Cross-range extent (m) & Detection span in azimuth — vessel lateral size proxy \\
Range (m) & Slant range — retained as a confound covariate, not a vessel property \\
\bottomrule
\end{tabular}
\end{center}

% ── Statistical evidence ─────────────────────────────────────────────────────
\section{Statistical Evidence for the Feature Ceiling}

Cohen's~$d$ (standardised mean difference) quantifies inter-class separation per
feature. Values $d < 0.5$ indicate distributions that overlap substantially;
$d \geq 1.0$ is the practical threshold for reliable single-feature discrimination.

\begin{center}\small
\begin{tabular}{lrrrrrr}
\toprule
Feature & Fis--Tug & Fis--Car & Fis--Tan & Tug--Car & Tug--Tan & \textbf{Car--Tan} \\
\midrule
Log peak RCS        & 0.47 & 1.29 & 1.80 & 0.84 & 1.25 & \textbf{0.20} \\
Log total RCS       & 0.87 & 1.59 & 2.52 & 0.82 & 1.49 & \textbf{0.33} \\
RCS concentration   & 0.98 & 0.51 & 1.32 & 0.13 & 0.25 & \textbf{0.26} \\
Aspect ratio        & 0.24 & 0.49 & 1.00 & 0.16 & 0.53 & \textbf{0.44} \\
Footprint (m$^2$)  & 0.92 & 1.36 & 2.09 & 0.52 & 1.17 & \textbf{0.44} \\
Down-range (m)      & 0.29 & 1.14 & 1.45 & 0.98 & 1.28 & \textbf{0.05} \\
Cross-range (m)     & 0.90 & 0.98 & 1.15 & 0.13 & 0.38 & \textbf{0.56} \\
\midrule
\textbf{Maximum}    & 0.98 & 1.59 & 2.52 & 0.98 & 1.49 & \textbf{0.56} \\
\bottomrule
\end{tabular}
\end{center}

Every class pair except Cargo--Tanker achieves $d \geq 0.98$ on at least one feature.
Cargo and Tanker reach a ceiling of $d = 0.56$ across all features combined.
Both vessel types operate at identical median ranges (30\,km), ruling out range bias.
Cargo and Tanker are both large steel-hulled vessels of overlapping physical size;
their amplitude-only radar signatures are near-indistinguishable by any combination
of the features above.

% ── Exhaustive feature engineering ────────────────────────────────────────────
\section{Exhaustive Feature Engineering: Ceiling Confirmed}

Three additional computational approaches were applied to verify the ceiling
is not a processing artefact and that no transformation of the existing
amplitude data recovers the missing discriminatory information.

\textbf{(1) Temporal fractal complexity of per-track RCS sequences.}
Higuchi fractal dimension, Petrosian FD, sample entropy, and Hurst exponent
were computed on the scan-to-scan RCS amplitude sequence for each vessel track.
These quantify whether the \emph{fluctuation pattern} across antenna rotations
— rather than the mean amplitude level — differs between vessel structural types.
Best Cargo--Tanker $d$: Hurst exponent on log-total-RCS $= 0.245$;
Higuchi FD $= 0.160$.
Neither approaches the $d = 0.56$ ceiling; the fluctuation pattern does not
encode structural differences accessible from the existing processed data.

\textbf{(2) Feature-space self-similarity (local intrinsic dimensionality).}
Local Intrinsic Dimensionality (LID) and multiscale neighbourhood density
were computed per detection across five feature subspaces to test whether
Cargo and Tanker occupy geometrically distinct regions of the feature space —
independent of their marginal amplitude distributions.

\begin{center}\small
\begin{tabular}{lcc}
\toprule
Subspace & Best Car--Tan $d$ & Feature type \\
\midrule
Amplitude only (3-D) & 0.393 & Neighbourhood density \\
RCS-5 features (5-D) & 0.333 & Neighbourhood density \\
Without range (7-D)  & 0.226 & Correlation dimension \\
Full feature set (8-D) & 0.194 & Neighbourhood density \\
Geometry only (4-D)  & 0.185 & Local intrinsic dimension \\
\midrule
\textbf{Amplitude feature ceiling} & \textbf{0.560} & Cross-range extent \\
\bottomrule
\end{tabular}
\end{center}

Cargo and Tanker detections are \textbf{geometrically co-located} in every
subspace tested: same neighbourhood density, same local dimensionality,
same correlation dimension. Restricting to the amplitude-only subspace
gives the highest self-similarity separation ($d = 0.393$), but still
30\% below the direct feature ceiling. Adding geometric dimensions
progressively dilutes rather than reveals structure.

\textbf{Conclusion.} The $d = 0.56$ ceiling is not a feature engineering
gap. No nonlinear transformation, fractal measure, or geometric analysis
of the existing amplitude-based features recovers additional discriminatory
information. The physically missing observable must originate from a
different property of the radar return — one not present in the
post-integration per-scan amplitude output of the current processing chain.

% ── The question ─────────────────────────────────────────────────────────────
\section{The Question Posed to the Expert}

\begin{quote}
\itshape
Given a fixed-hardware, amplitude-only, single-polarisation marine surveillance
radar: does the electromagnetic return from a large steel vessel at 20--50\,km
encode physical information that distinguishes a structurally complex vessel
(container ship: heterogeneous deck cargo, cranes, hatches) from a structurally
simpler vessel (tanker: smooth low-profile hull and deck), and if so, can this
information be extracted through receiver tuning or processing parameter changes
alone — without any hardware modification?
\end{quote}

Specific aspects where expert insight is sought:

\begin{enumerate}[noitemsep, leftmargin=*, label=\arabic*.]

  \item \textbf{Scan-to-scan amplitude fluctuation statistics.}
  The vessel RCS changes from one antenna rotation to the next. A container ship's
  structurally complex deck may produce a different fluctuation regime (Swerling
  model, variance, autocorrelation length) than a smooth tanker hull.
  \textit{Is this distinction physically expected? How many consecutive scans are
  needed to estimate it reliably, and which receiver integration setting best
  preserves this information?}

  \item \textbf{Within-dwell amplitude distribution.}
  Multiple pulses illuminate the vessel per beam dwell. The per-pulse amplitude
  sequence before integration may encode structural complexity.
  \textit{Is this within-dwell distribution (skewness, pulse-to-pulse variance)
  expected to differ between vessel structural types? Is it accessible through
  receiver output logging without hardware change?}

  \item \textbf{Doppler content and clutter filter tuning.}
  The existing clutter filter suppresses slow targets including vessel Doppler.
  \textit{If the filter cutoff is lowered (a tunable parameter), is bulk vessel
  Doppler or sub-resolution scatterer motion (propeller, rotating machinery)
  expected to appear in the output? Would that motion signature differ between
  a diesel-direct-drive cargo vessel and a slow-speed tanker drive system,
  and at what filter setting would it become accessible?}

  \item \textbf{Clutter interaction geometry.}
  The vessel modifies the surrounding sea clutter field (shadow, bow-wave, wake).
  This is present in the radar picture but not currently extracted.
  \textit{Is this interaction geometry expected to differ between vessel types?
  What detection threshold or CFAR guard-band adjustment would expose it?}

  \item \textbf{Any other receiver-accessible physical observable.}
  \textit{Is there a property of the radar-target interaction — not listed above
  — that is physically expected to correlate with vessel structural type and
  is accessible through receiver settings or processing changes on this
  fixed-hardware system?}

\end{enumerate}

For any property identified, the author requests an assessment of: (a)~its physical
origin, (b)~the specific receiver or processing parameter that controls its
observability, and (c)~a rough expectation of the inter-class separability relative
to the $d = 0.56$ ceiling currently achieved.

\vspace{6pt}
\noindent\rule{\linewidth}{0.4pt}\\[2pt]
\noindent\small\textit{This document is shared in strict confidence.
The associated classification methodology and findings are the intellectual
property of Vijay Kumar V and have not been publicly disclosed.
Please do not circulate without the author's written permission.}

\end{document}
